\section{单位元的添加}

设 A 是一个 k-代数。在 III, §1, No.~2, 第 431 页中，我们定义了由 A 通过添加单位元素 e 而得到的有单位元代数 \(\widetilde{A}\)。我们把 A 与 \(\widetilde{A}\) 的一个双边理想等同起来。k-模 \(\widetilde{A}\) 是子模 ke 与 A 的直和。

命题 3. —— a) 设 \(\widetilde{\mathfrak{a}}\) 是 \(\widetilde{\mathrm{A}}\) 的一个左理想，使得 \(\widetilde{\mathrm{A}} = \widetilde{\mathfrak{a}} + \mathrm{A}\)。我们置 \(\mathfrak{a} = \widetilde{\mathfrak{a}} \cap \mathrm{A}\)。存在一个元素 \(u\)（\(\mathrm{A}\) 的元素），使得 \(u - e\) 属于 \(\widetilde{\mathfrak{a}}\)。若 \(u\) 是这样一个元素，则它是 \(\mathrm{A}\) 的模 \(\mathfrak{a}\) 的右单位元，并且我们有 \(\widetilde{\mathfrak{a}} = \mathfrak{a} + k(u - e)\)；特别地，理想 \(\mathfrak{a}\) 是正则的。

\begin{enumerate}
\def\labelenumi{\alph{enumi})}
\setcounter{enumi}{1}
\item
  反之，设 \(\mathfrak{a}\) 是 A 的一个正则左理想，\(u\) 是 A 的模 \(\mathfrak{a}\) 的一个右单位元。置 \(\widetilde{\mathfrak{a}} = \mathfrak{a} + k(u - e)\)。则 \(\widetilde{\mathfrak{a}}\) 是 \(\widetilde{\mathrm{A}}\) 的一个左理想，使得 \(\widetilde{\mathrm{A}} = \widetilde{\mathfrak{a}} + \mathrm{A}\)，并且我们有 \(\mathfrak{a} = \widetilde{\mathfrak{a}} \cap \mathrm{A}\)。
\item
  若 \(k\) 是一个体，则条件 \(\widetilde{\mathrm{A}} = \widetilde{\mathfrak{a}} + \mathrm{A}\) 等价于说 \(\widetilde{\mathfrak{a}}\) 不包含于 \(\mathrm{A}\) 之中。
\end{enumerate}

在 a) 的假设下，元素 \(e\)（\(\widetilde{\mathbf{A}}\) 的元素）可以写成 \((e - u) + u\)，其中 \(u \in \mathbf{A}\) 且 \(u - e \in \widetilde{\mathfrak{a}}\)。设 \(x = \lambda e + a\) 是 \(\widetilde{\mathbf{A}}\) 的一个元素，其中 \(\lambda \in k\) 且 \(a \in \mathbf{A}\)；若 \(x\) 属于 \(\widetilde{\mathfrak{a}}\)，则元素 \(y = a + \lambda u = x + \lambda (u - e)\)（\(\mathbf{A}\) 的元素）属于 \(\mathfrak{a}\)，并且我们有 \(x = y - \lambda (u - e)\)；这就证明了等式 \(\widetilde{\mathfrak{a}} = \mathfrak{a} + k(u - e)\)。此外，对每个 \(a\)（\(\mathbf{A}\) 中的元素），元素 \(au - a = a(u - e)\)（\(\mathbf{A}\) 的元素）属于 \(\mathfrak{a}\)，因此 \(u\) 是模 \(\mathfrak{a}\) 的右单位元。这就证明了断言 a)。

设 \(\mathfrak{a}\) 与 \(u\) 如 b) 中所述。由于 \(\widetilde{\mathrm{A}}\) 是 \(\mathrm{A}\) 与 \(k(u - e)\) 的直和，我们有 \(\widetilde{\mathfrak{a}} + \mathrm{A} = \widetilde{\mathrm{A}}\) 与 \(\widetilde{\mathfrak{a}} \cap \mathrm{A} = \mathfrak{a}\)。\(\widetilde{\mathfrak{a}}\) 的每个元素都形如 \(x + \lambda (u - e)\)，其中 \(x \in \mathfrak{a}\) 且 \(\lambda \in k\)。对每个 \(a \in \mathrm{A}\)，我们有

\[
a (x + \lambda (u - e)) = a x + \lambda (a u - a).
\]

由于 u 是 A 的模 \(\mathfrak{a}\) 的右单位元，和 \(ax + \lambda(au - a)\) 属于 \(\mathfrak{a}\)，从而 \(\widetilde{\mathfrak{a}}\) 是 A 的一个左理想。这就证明了 b)。

最后，若 \(k\) 是一个体，则 \(A\) 是 \(\widetilde{A}\) 中的一个超平面，而 \(\widetilde{A}\) 是 \(\widetilde{\mathfrak{a}}\) 与 \(A\) 之和，当且仅当 \(\widetilde{\mathfrak{a}}\) 不包含于 \(A\) 之中。

推论. —— A 的正则左理想是形如 \(A \cap \widetilde{\mathfrak{a}}\) 的理想，其中 \(\widetilde{\mathfrak{a}}\) 是 \(\widetilde{A}\) 的一个左理想，使得 \(\widetilde{A} = \widetilde{\mathfrak{a}} + A\)。

命题 4. —— a) 设 \(\mathfrak{a}\) 是 A 的一个正则极大左理想。存在唯一的左理想 \(\widetilde{\mathfrak{a}}\)（\(\widetilde{\mathrm{A}}\) 的左理想），使得 \(\widetilde{\mathrm{A}} = \widetilde{\mathfrak{a}} + \mathrm{A}\) 且 \(\mathfrak{a} = \widetilde{\mathfrak{a}} \cap \mathrm{A}\)。这个理想是极大的且不包含 A。

\begin{enumerate}
\def\labelenumi{\alph{enumi})}
\setcounter{enumi}{1}
\tightlist
\item
  映射 \(\widetilde{\mathfrak{a}}\mapsto \widetilde{\mathfrak{a}}\cap \mathrm{A}\) 是从 \(\widetilde{\mathrm{A}}\) 的不包含 A 的极大左理想所成之集到 A 的正则极大左理想所成之集上的一个双射。
\end{enumerate}

设 \(\mathfrak{a}\) 是 A 的一个正则极大左理想。由命题 3, a)，左理想 \(\mathfrak{b}\)（\(\widetilde{\mathrm{A}}\) 的左理想，满足 \(\mathfrak{b} + \mathrm{A} = \widetilde{\mathrm{A}}\) 与 \(\mathfrak{b} \cap \mathrm{A} = \mathfrak{a}\)）就是诸理想 \(\mathfrak{a} + k(u - e)\)，其中 \(u\) 是模 \(\mathfrak{a}\) 的右单位元。因此，为证明 \(\widetilde{\mathfrak{a}}\) 的唯一性，只需证明 A 的两个右单位元 \(u\) 与 \(u'\)（模 \(\mathfrak{a}\) 的右单位元）模 \(\mathfrak{a}\) 同余即可。用反证法推理，假设 \(u - u'\) 不属于 \(\mathfrak{a}\)。公式 \(x(u - u') = (xu - x) - (xu' - x)\) 表明我们有 \(\mathrm{A}(u - u') \subset \mathfrak{a}\)；由此推出 \(\mathfrak{a} + k(u - u')\) 是 A 的一个包含 \(\mathfrak{a}\) 且异于 \(\mathfrak{a}\) 的左理想。由于 \(\mathfrak{a}\) 是极大的，我们因此有 \(\mathfrak{a} + k(u - u') = \mathrm{A}\)，从而 \(\mathrm{AA} \subset \mathfrak{a}\)。对每个 \(x \in \mathrm{A}\)，我们有 \(x \equiv xu (\mathrm{mod} \mathfrak{a})\)，因此由上所述 \(x \in \mathfrak{a}\)，这与假设 \(\mathfrak{a} \neq \mathrm{A}\) 矛盾。

于是存在唯一的左理想 \(\widetilde{\mathfrak{a}}\)（\(\widetilde{\mathrm{A}}\) 的左理想），使得 \(\widetilde{\mathrm{A}} = \widetilde{\mathfrak{a}} + \mathrm{A}\) 且 \(\mathfrak{a} = \widetilde{\mathfrak{a}} \cap \mathrm{A}\)。设 \(\mathfrak{b}\) 是 \(\widetilde{\mathrm{A}}\) 的一个包含 \(\widetilde{\mathfrak{a}}\) 的真左理想。则 \(\mathfrak{b} \cap \mathrm{A}\) 是 \(\mathrm{A}\) 的一个包含 \(\mathfrak{a}\) 的真左理想。由于 \(\mathfrak{a}\) 是极大的，它因此等于 \(\mathfrak{a}\)，从而由 VIII, 第 4 页的引理 2 推出 \(\mathfrak{b} = \widetilde{\mathfrak{a}}\)。这就证明了 \(\widetilde{\mathfrak{a}}\) 是 \(\widetilde{\mathrm{A}}\) 的一个极大理想；对这样一个理想，条件 \(\widetilde{\mathrm{A}} = \widetilde{\mathfrak{a}} + \mathrm{A}\) 意味着 \(\widetilde{\mathfrak{a}}\) 不包含 A。

剩下的要证明：若 \(\widetilde{\mathfrak{a}}\) 是 \(\tilde{A}\) 的一个不包含 A 的极大左理想，则 A 的左理想 \(\mathfrak{a} = A \cap \widetilde{\mathfrak{a}}\) 是极大的。设 \(\mathfrak{b}\) 是 A 的一个包含 \(\mathfrak{a}\) 的真左理想。设 u 是模 \(\mathfrak{a}\) 的一个右单位元，使得 \(\widetilde{\mathfrak{a}} = \mathfrak{a} + k(u - e)\)（命题 3）。它也是模 \(\mathfrak{b}\) 的右单位元；用 \(\widetilde{\mathfrak{b}}\) 表示理想 \(\mathfrak{b} + k(u - e)\)（\(\tilde{A}\) 的理想）。由命题 3, b)，我们有 \(\mathfrak{b} = A \cap \widetilde{\mathfrak{b}}\)。理想 \(\widetilde{\mathfrak{a}}\)（等于 \(\mathfrak{a} + k(u - e)\) 的理想）包含于 \(\widetilde{\mathfrak{b}}\)，从而等于 \(\widetilde{\mathfrak{b}}\)，因为 \(\widetilde{\mathfrak{b}}\) 异于 \(\tilde{A}\) 而 \(\widetilde{\mathfrak{a}}\) 极大。从而我们有 \(\mathfrak{a} = \mathfrak{b}\)，故 \(\mathfrak{a}\) 是极大的。

我们把命题 3 与命题 4 向右理想的翻译留给读者。

